% ══════════════════════════════════════════════════════════════════════
% RELATÓRIO CURTO: Indicadores dos municípios de São Paulo
% ══════════════════════════════════════════════════════════════════════
% Feito de estilo/modelo_relatorio.tex (o kit do portfólio). O bloco entre
% "PREÂMBULO DO KIT" e "FIM DO PREÂMBULO DO KIT" é o do kit, sem mudança.
%
% Nenhum número é digitado aqui: cada \Num... vem de relatorio/numeros.tex,
% e as linhas de tabela e os trechos de código vêm de relatorio/gerados/,
% todos escritos por scripts/numeros.py. Registro em NUMEROS_CANONICOS.md.
%
% Compilar (da pasta relatorio/): pdflatex relatorio.tex, duas passadas.
% Ou, da raiz: python scripts/numeros.py && python scripts/verificar_textos.py
% (o verificador compila e grava relatorio/relatorio.pdf).

% ══════════════════════════════════════════════════════════════════════
% PREÂMBULO DO KIT: daqui até "FIM DO PREÂMBULO DO KIT", não edite
% ══════════════════════════════════════════════════════════════════════
\documentclass[11pt,a4paper]{article}

\usepackage[T1]{fontenc}
\usepackage[utf8]{inputenc}
\usepackage[brazilian]{babel}

% A mancha: 15 cm de linha, cerca de 3,2 alfabetos de Libertine 11pt. As
% figuras NÃO vivem nesta medida: ver \figwidth, mais abaixo.
\usepackage[top=2cm, bottom=2cm, left=3cm, right=3cm]{geometry}
\usepackage{graphicx}
\usepackage{subcaption}
\usepackage{booktabs}
\usepackage{array}
\usepackage{longtable}
\usepackage{float}
\usepackage{amsmath}
\usepackage{caption}
% A opção `table` liga o colortbl, que é o que dá a zebra (\rowcolors).
\usepackage[table]{xcolor}
\usepackage{enumitem}
% Capa: a curva em TikZ e os ícones vetoriais do Font Awesome.
\usepackage{tikz}
\usepackage{fontawesome5}

% ── AS FONTES: serifa no corpo, Lato em tabela e figura ───────────────
% `lato` SEM a opção `default`: o corpo não muda e o comando \lato fica
% disponível. A `libertine` entra DEPOIS da `lato`, porque as duas mexem
% em \sfdefault e a última vence. A matemática acompanha a Libertine.
\usepackage{lato}
\usepackage{libertine}
\usepackage{libertinust1math}

% ── AS CORES: ColorBrewer Set1 e a escala tirada do #0A0310 ───────────
% Os mesmos hexadecimais do estilo_portfolio.py e do theme_portfolio.R.
\definecolor{azul}{HTML}{377EB8}       % destaque principal
\definecolor{vermelho}{HTML}{E41A1C}   % destaque secundário
\definecolor{verde}{HTML}{4DAF4A}
\definecolor{roxo}{HTML}{984EA3}
\definecolor{laranja}{HTML}{FF7F00}
\definecolor{marrom}{HTML}{A65628}
\definecolor{rosa}{HTML}{F781BF}
\definecolor{tinta}{HTML}{0A0310}      % o quase-preto do kit
\definecolor{tinta2}{HTML}{403A45}     % GRAY_70: texto secundário
\definecolor{tinta3}{HTML}{5D5961}     % GRAY_60: rodapé, nota
% Versões de TEXTO das cores claras, escurecidas até 4,5:1 contra o branco
% por cor_de_texto() do estilo_portfolio.py. Vermelho e roxo já passam.
\definecolor{azulTexto}{HTML}{3579B1}
\definecolor{verdeTexto}{HTML}{3C823B}

% ── A ZEBRA, aplicada a TODA tabela sem editar tabela nenhuma ─────────
% A alternância é ancorada no PRIMEIRO \midrule de cada tabular (o fim do
% cabeçalho), e não no número da linha: assim a primeira linha de dado
% sai sempre branca, com cabeçalho de uma ou de duas linhas.
%
% ⚠ A ordem {zebraclara}{white} é contraintuitiva e está certa: o
% contador sobe ANTES de a cor ser escolhida. Invertê-la inverte a zebra.
% ⚠ O reset vai dentro de \noalign, senão "Misplaced \noalign".
\usepackage{etoolbox}
\definecolor{zebraclara}{HTML}{F4F4F4}
\newif\ifzebraprimeiromidrule
\newcommand{\zebratabela}{%
  \lato
  \global\zebraprimeiromidruletrue
  \rowcolors{1}{zebraclara}{white}%
}
\AtBeginEnvironment{tabular}{\zebratabela}
\AtBeginEnvironment{longtable}{\zebratabela}
\pretocmd{\midrule}{%
  \noalign{%
    \ifzebraprimeiromidrule
      \global\zebraprimeiromidrulefalse
      \global\rownum=0\relax
    \fi}%
}{}{\typeout{ATENCAO: o patch de \string\midrule\space para a zebra FALHOU}}

% ── A BOLINHA DE ESTADO, para coluna sim/não de tabela ────────────────
% É reforço, e não substituição: o texto fica ao lado. Os três tons têm o
% mesmo contraste (2,3:1), calculado no projeto de origem. Glifo, e não
% TikZ: \tikz dentro de \caption derruba a compilação.
\definecolor{bolinhaSim}{HTML}{6DBD6B}
\definecolor{bolinhaNao}{HTML}{F29192}
\definecolor{bolinhaEmpate}{HTML}{B1AE28}
\DeclareRobustCommand{\bolinha}[1]{\textcolor{#1}{%
  \raisebox{-0.06ex}{\scalebox{2}{$\bullet$}}}}
\DeclareRobustCommand{\marcasim}{\bolinha{bolinhaSim}\,sim}
\DeclareRobustCommand{\marcanao}{\bolinha{bolinhaNao}\,não}
\DeclareRobustCommand{\marcaempate}{\bolinha{bolinhaEmpate}\,empate}

% ── A CAIXA DE DESTAQUE ───────────────────────────────────────────────
% Para o achado principal de uma seção, uma por seção no máximo. Fio azul
% à esquerda, fundo da zebra, título em Lato. Quebra entre páginas.
%   \begin{destaque}[Título opcional] texto \end{destaque}
\usepackage[most]{tcolorbox}
\newtcolorbox{destaque}[1][Em uma frase]{%
  enhanced, breakable,
  colback=zebraclara, colframe=azul,
  boxrule=0pt, leftrule=2.5pt, arc=0pt, outer arc=0pt,
  left=10pt, right=10pt, top=8pt, bottom=8pt,
  before skip=12pt, after skip=12pt,
  % ⚠ A cor do título é `coltitle`: um \color antes de \tcbtitle não vale,
  % porque o tcolorbox reaplica coltitle (branco por padrão) por dentro.
  title={#1}, detach title, coltitle=azulTexto,
  before upper={{\lato\bfseries\tcbtitle}\par\smallskip},
}

\usepackage{needspace}
\usepackage[hidelinks]{hyperref}
\setcounter{tocdepth}{2}
\hypersetup{bookmarksdepth=3, bookmarksnumbered=true, bookmarksopen=false}

\captionsetup{font=small, labelfont=bf}
\captionsetup[sub]{font=footnotesize}

\graphicspath{{figuras/}{exemplos/}}

% Paginação: altura natural, sem esticar o espaço vertical; sem órfãs nem
% viúvas; folga mínima de composição em vez de reescrever prosa.
\raggedbottom
\widowpenalty=10000
\clubpenalty=10000
\emergencystretch=1em

% ── \figwidth: a figura é mais larga que a prosa ──────────────────────
% A figura vive na medida de 1 cm de margem no papel (19 cm em A4), e não
% na mancha de 15 cm. Definida pelo papel, e não por número copiado.
\newlength{\figwidth}
\setlength{\figwidth}{\dimexpr\paperwidth-2cm\relax}

% \imagemplena[fração]{arquivo}: a imagem centrada sobre a mancha. A caixa
% de largura declarada impede o transbordo só para a direita e não conta
% como Overfull \hbox.
\newcommand{\imagemplena}[2][1]{%
  \makebox[\linewidth][c]{\includegraphics[width=#1\figwidth]{#2}}}

% tabelaplena: a tabela que não cabe na mancha ganha a medida da figura.
% Só para tabular (caixa não quebra entre páginas).
\newsavebox{\tabelacaixa}
\newenvironment{tabelaplena}
  {\begin{lrbox}{\tabelacaixa}}
  {\end{lrbox}\makebox[\linewidth][c]{\usebox{\tabelacaixa}}}

% \subfig[largura]{arquivo}{legenda}: um painel por linha, na medida da figura.
\newcommand{\subfig}[3][0.62\figwidth]{%
  \makebox[\linewidth][c]{%
  \begin{subfigure}[t]{#1}
    \centering
    \includegraphics[width=\linewidth]{#2}
    \caption{#3}
  \end{subfigure}}\par\medskip}
% \linhasde{arquivo}: as linhas de uma tabela, escritas por script.
% ⚠ Não use \input dentro de tabular: ele não é expansível, e o
% \bottomrule seguinte derruba a compilação com "Misplaced \noalign". O
% \@@input primitivo é expansível, e o espaço depois de #1 encerra o nome.
\makeatletter
\newcommand{\linhasde}[1]{\@@input #1 }
\makeatother
% ══════════════════════════════════════════════════════════════════════
% FIM DO PREÂMBULO DO KIT
% ══════════════════════════════════════════════════════════════════════

% ── AJUSTES DO PROJETO (fora do bloco do kit) ─────────────────────────
% As figuras e as capturas ficam fora de relatorio/.
\graphicspath{{../graficos/}{../powerbi/capturas/}}

% A captura do Power BI entra se o arquivo existir; senão, uma caixa diz
% que ela falta, e o PDF compila do mesmo jeito. Uso:
%   \def\capturaArquivo{../powerbi/capturas/arquivo.png}\captura
\newtcolorbox{capturapendente}{%
  colback=zebraclara, colframe=tinta, boxrule=0.4pt, arc=0pt,
  width=0.7\figwidth, height=4cm, valign=center, halign=center, center}
\newcommand{\captura}{%
  \IfFileExists{\capturaArquivo}%
    {\imagemplena[0.86]{\capturaArquivo}}%
    {\begin{capturapendente}
       {\lato\color{tinta2}Captura pendente: será exportada do Power BI Desktop\par
        \smallskip\footnotesize\texttt{\capturaArquivo}}
     \end{capturapendente}}}

% codigo: trecho de código (verbatim) que começa na margem do texto e pode
% avançar sobre a margem direita, até a medida da figura.
\newsavebox{\codigocaixa}
\newenvironment{codigo}
  {\par\smallskip\begin{lrbox}{\codigocaixa}\begin{minipage}{1.2\linewidth}\footnotesize}
  {\end{minipage}\end{lrbox}\noindent\makebox[\linewidth][l]{\usebox{\codigocaixa}}\par\smallskip}

% ── DADOS DA CAPA: preencha por relatório ─────────────────────────────
\newcommand{\RotuloCapa}{RELATÓRIO TÉCNICO \textperiodcentered\ PROJETO DE PORTFÓLIO}
\newcommand{\TituloRelatorio}{Indicadores dos municípios de São Paulo}
\newcommand{\SubtituloRelatorio}{Modelo estrela no Power BI, planilha com PROCX e previsão da receita testada contra o passado}
\newcommand{\DelimitacaoRelatorio}{\NumMunicipios{} municípios paulistas, \NumIndicadores{} Indicadores de dado público \\ Período: \NumAnoInicial{} a \NumAnoFinal; previsão até \NumAnoPrevFinal}
\newcommand{\AutorNome}{João Eduardo Pastori Garcia}
\newcommand{\AutorPapel}{Graduando em Estatística, Unicamp}
\newcommand{\AutorLinkedIn}{https://www.linkedin.com/in/joao-eduardo-garcia/}
\newcommand{\LocalCapa}{São Paulo -- SP}
\newcommand{\DataCapa}{\NumMesCapa{} de \NumAnoCapa}

% ── OS NÚMEROS DO RELATÓRIO: gerados por script, nunca digitados ──────
% GERADO por scripts/numeros.py. Não edite à mão: rode o script de novo.
% Registro de cada número em NUMEROS_CANONICOS.md.
\newcommand{\NumMunicipios}{645}
\newcommand{\NumIndicadores}{8}
\newcommand{\NumIndicadoresComparaveis}{7}
\newcommand{\NumAnoInicial}{2015}
\newcommand{\NumAnoFinal}{2025}
\newcommand{\NumLinhasFato}{50.413}
\newcommand{\NumFaixas}{7}
\newcommand{\NumMunicipioAnos}{7.095}
\newcommand{\NumExclusoes}{27}
\newcommand{\NumMunicipiosComExclusao}{15}
\newcommand{\NumAnoPrevInicial}{2026}
\newcommand{\NumAnoPrevFinal}{2028}
\newcommand{\NumMunicipiosComPrevisao}{642}
\newcommand{\NumAnoBase}{2025}
\newcommand{\NumAnoFinalPib}{2023}
\newcommand{\NumNivelMenor}{80}
\newcommand{\NumNivelMaior}{95}
\newcommand{\NumPaginasPainel}{4}
\newcommand{\NumDataConsultaFontes}{14/10/2026}
\newcommand{\NumTabelasFato}{3}
\newcommand{\NumDimensoes}{3}
\newcommand{\NumRelacoes}{7}
\newcommand{\NumMedidas}{39}
\newcommand{\NumCelulasConferencia}{174}
\newcommand{\NumMunicipiosConferencia}{3}
\newcommand{\NumCelulasBatem}{174}
\newcommand{\NumCelulasCompletas}{151.239}
\newcommand{\NumDivergenciasCompletas}{0}
\newcommand{\NumExemplo}{Ibirá}
\newcommand{\NumExemploCandidatos}{521}
\newcommand{\NumExemploFaixa}{de 10 a 20 mil}
\newcommand{\NumExemploNaFaixa}{120}
\newcommand{\NumExemploPopulacao}{11.690}
\newcommand{\NumFigUmVerbo}{gasta}
\newcommand{\NumFigUmObjeto}{com saúde por habitante}
\newcommand{\NumFigUmDif}{18,7}
\newcommand{\NumFigUmSentido}{mais}
\newcommand{\NumFigUmAno}{2025}
\newcommand{\NumDifReceita}{$-$4,9}
\newcommand{\NumDifSaude}{+18,7}
\newcommand{\NumDifEducacao}{$-$7,1}
\newcommand{\NumDifDependencia}{$-$1,0}
\newcommand{\NumFigDoisAcima}{3}
\newcommand{\NumFigDoisExtremo}{Ideb dos anos iniciais, rede municipal}
\newcommand{\NumFigDoisExtremoPos}{92}
\newcommand{\NumExemploReceitaUltima}{76,7}
\newcommand{\NumExemploPrevista}{88,6}
\newcommand{\NumExemploInfOitenta}{74,9}
\newcommand{\NumExemploSupOitenta}{114,6}
\newcommand{\NumExemploInfNoventaCinco}{67,6}
\newcommand{\NumExemploSupNoventaCinco}{128,2}
\newcommand{\NumErroEscolhidoUm}{6,6}
\newcommand{\NumErroEscolhidoDois}{9,4}
\newcommand{\NumErroEscolhidoTres}{12,5}
\newcommand{\NumViesEscolhidoTres}{+5,2}
\newcommand{\NumErroIngenuoUm}{7,0}
\newcommand{\NumErroIngenuoDois}{12,7}
\newcommand{\NumErroIngenuoTres}{18,4}
\newcommand{\NumViesIngenuoTres}{+18,2}
\newcommand{\NumErroRetaUm}{6,8}
\newcommand{\NumErroRetaDois}{10,0}
\newcommand{\NumErroRetaTres}{13,3}
\newcommand{\NumViesRetaTres}{+9,1}
\newcommand{\NumOrigensTeste}{4}
\newcommand{\NumOrigemInicial}{2019}
\newcommand{\NumOrigemFinal}{2022}
\newcommand{\NumCobOitentaSem}{71}
\newcommand{\NumCobOitentaCom}{81}
\newcommand{\NumCobOitentaPior}{62}
\newcommand{\NumCobOitentaPiorOrigem}{2022}
\newcommand{\NumEscalaOitenta}{1,25}
\newcommand{\NumCobNoventaCincoSem}{92}
\newcommand{\NumCobNoventaCincoCom}{95}
\newcommand{\NumCobNoventaCincoPior}{91}
\newcommand{\NumCobNoventaCincoPiorOrigem}{2022}
\newcommand{\NumEscalaNoventaCinco}{1,15}
\newcommand{\NumCelulasExcel}{256}
\newcommand{\NumCelulasExcelBatem}{256}
\newcommand{\NumMunicipiosExcel}{4}


\title{\TituloRelatorio}

\begin{document}

% ══════════════════════════════ CAPA ══════════════════════════════════
\begin{titlepage}
  \centering
  \vspace*{0.5cm}

  {\color{tinta}\rule{\linewidth}{1.5pt}} \\[0.4cm]
  {\small\bfseries\color{tinta2} \RotuloCapa} \\[0.2cm]
  {\color{tinta}\rule{\linewidth}{0.5pt}}

  \vspace{1.6cm}

  \begin{tikzpicture}
    \draw[line width=1.6pt, azul, smooth]
      plot coordinates {(0.2,0.35) (1.7,0.62) (3.2,0.48) (4.7,0.95)
                        (6.2,1.32) (7.7,1.12) (9.2,1.58) (10.7,1.42)};
    \foreach \x/\c in {0.2/verde, 2.9/vermelho, 5.6/azul, 8.3/roxo}
      \fill[\c] (\x,-0.28) rectangle ++(2.4,0.07);
  \end{tikzpicture}

  \vspace{1.2cm}

  \begin{minipage}{0.94\textwidth}
    \centering
    {\Huge\bfseries\color{azul} \TituloRelatorio \par}
    \vspace{0.5cm}
    {\LARGE\bfseries\color{vermelho} \SubtituloRelatorio \par}
  \end{minipage}

  \vspace{1.1cm}

  % Um ícone por eixo: a base, o modelo estrela, a planilha e a previsão.
  {\Large
   \textcolor{azul}{\faDatabase}\hspace{1.1cm}
   \textcolor{vermelho}{\faProjectDiagram}\hspace{1.1cm}
   \textcolor{azul}{\faTable}\hspace{1.1cm}
   \textcolor{vermelho}{\faChartLine}\par}

  \vspace{1.0cm}

  {\large\itshape\color{tinta2} \DelimitacaoRelatorio \par}

  \vspace{2.2cm}

  {\large\bfseries\color{azul} Autoria} \\[0.6cm]

  \begin{minipage}{0.85\textwidth}
    \centering
    {\large\bfseries \AutorNome} \\
    {\normalsize\color{tinta2} \AutorPapel} \\[0.15cm]
    {\footnotesize
     \href{\AutorLinkedIn}{\color{vermelho}\textbf{LinkedIn}}}
  \end{minipage}

  \vfill

  {\color{tinta2}\rule{4cm}{0.5pt}} \\[0.3cm]
  {\small\color{tinta2} \LocalCapa \\ \bfseries \DataCapa}
\end{titlepage}

\tableofcontents
\clearpage

% ══════════════════════════════════════════════════════════════════════
\section{Onde o município está, e onde estão os parecidos}
\label{sec:problema}

Quem acompanha as contas e os serviços de uma prefeitura tem o dado, mas
não tem como responder, sem uma semana de planilha, a três perguntas:
\textbf{onde o meu município está}, \textbf{onde estão os municípios
parecidos com ele} e \textbf{para onde a receita aponta}. O dado existe,
aberto, em quatro lugares (Tesouro Nacional, IBGE, INEP e DataSUS), em
formatos e anos diferentes e com erros de declaração.

O projeto responde às três perguntas para os \NumMunicipios{} municípios
de São Paulo, de \NumAnoInicial{} a \NumAnoFinal, com \NumIndicadores{}
Indicadores definidos por escrito (fórmula, fonte, periodicidade e meta).
Os Municípios parecidos são os da mesma Faixa de porte do Censo 2022, e a
comparação mostra a Mediana da faixa, os quartis e a Posição na faixa,
nunca um ranking por nome. As entregas são um painel no Power BI sobre um
modelo estrela, uma planilha de Excel com PROCX e tabela dinâmica e uma
Previsão da receita testada contra o passado.

O município de exemplo de todas as figuras é \NumExemplo{}, com
\NumExemploPopulacao{} habitantes no Censo 2022, sorteado com semente fixa
entre os \NumExemploCandidatos{} municípios com série completa: ninguém o
escolheu pelo que ele mostra. A Figura~\ref{fig:painel} responde à
primeira pergunta para ele; a Figura~\ref{fig:faixa}, à segunda.

\begin{figure}[H]
  \centering
  \imagemplena[0.86]{01-painel-municipio.png}
  \caption{\NumExemplo{} \NumFigUmVerbo{} \NumFigUmDif\% \NumFigUmSentido{}
    \NumFigUmObjeto{} que a mediana dos \NumExemploNaFaixa{} municípios
    \NumExemploFaixa{} habitantes, em \NumFigUmAno. Valor do município em
    vermelho, Mediana da faixa em cinza e, na sombra, do primeiro ao terceiro
    quartil dos Municípios parecidos, em reais de \NumAnoBaseReais.}
  \label{fig:painel}
\end{figure}

\begin{figure}[H]
  \centering
  \imagemplena[0.86]{02-municipio-na-faixa.png}
  \caption{\NumExemplo{} fica acima da Mediana da faixa em
    \NumFigDoisAcima{} de \NumIndicadoresComparaveis{} Indicadores. O mais
    perto de uma ponta da faixa é o \NumFigDoisExtremo, com Posição na faixa
    de \NumFigDoisExtremoPos\%. Cada caixa vai do primeiro ao terceiro quartil
    dos Municípios parecidos, no último ano com dado de cada Indicador.}
  \label{fig:faixa}
\end{figure}

% ══════════════════════════════════════════════════════════════════════
\section{Saíram \NumExclusoes{} valores declarados, de \NumMunicipioAnos{} declarações}
\label{sec:dados}

As finanças vêm da Declaração de Contas Anuais (DCA, anexos I-C e I-E) do
Siconfi, do Tesouro Nacional; a população, o PIB dos municípios e o IPCA,
do IBGE; o Ideb, do INEP; e os nascidos vivos por número de consultas de
pré-natal, do SINASC (DataSUS). Consultas SQL comentadas montam uma base
SQLite com uma tabela fato longa, de \NumLinhasFato{} linhas (município,
ano e Indicador com dado). Os valores em reais estão a preços de
\NumAnoBaseReais, pela média anual do IPCA.

A DCA é autodeclarada, e alguns municípios declararam, em certos anos,
valores várias vezes menores que os seus próprios dos outros anos. A
\textbf{regra de qualidade da declaração} é igual para todos e compara o
município com ele mesmo, e não com a faixa: valor zero ou negativo sai, e
o valor de um ano sai quando passa de \NumRegraFator{} vezes a média dos
\emph{outros} anos do mesmo município, para cima ou para baixo. Um
município de royalties, com receita alta em todos os anos, fica. O que sai
fica listado, com o motivo, num arquivo próprio; município que não
entregou a declaração fica sem valor, e nada é imputado.

\begin{destaque}[A regra da declaração, em números]
  Saíram \NumExclusoes{} valores, de \NumMunicipiosComExclusao{}
  municípios, entre \NumMunicipioAnos{} declarações possíveis
  (\NumMunicipios{} municípios em cada ano da série). A regra corta o erro
  de declaração sem tirar o município que é diferente dos outros de
  verdade.
\end{destaque}

Cada Indicador foi refeito em pandas, sem olhar o SQL, e comparado linha
a linha com a consulta (Seção~\ref{sec:conferencia}).

% ══════════════════════════════════════════════════════════════════════
\section{Um modelo estrela com \NumTabelasFato{} fatos e \NumDimensoes{} dimensões}
\label{sec:modelo}

O modelo semântico do Power BI está escrito em TMDL, texto que se lê e se
revisa no próprio repositório, linha a linha. As consultas do Power Query
leem os CSV tratados por um parâmetro de pasta. As tabelas fato ficam no
centro; as dimensões, em volta, ligadas por \NumRelacoes{} relações
um-para-muitos, todas de direção única, com as chaves ocultas
(Tabela~\ref{tab:modelo} e Figura~\ref{fig:modelo}).

\begin{table}[H]
  \centering
  \caption{As tabelas do modelo estrela, lidas do TMDL. A coluna da direita
    lista as tabelas com que cada uma se relaciona.}
  \label{tab:modelo}
  \small
  \begin{tabelaplena}
  \begin{tabular}{@{}lllrp{6cm}@{}}
    \toprule
    Tabela & Papel & Grão & Linhas & Liga-se a \\
    \midrule
    \linhasde{gerados/tabelas_modelo.tex}
    \bottomrule
  \end{tabular}
  \end{tabelaplena}
\end{table}

\textbf{Por que a referência da faixa é uma segunda fato.} A mediana e os
quartis vêm prontos, por Faixa de porte, ano e Indicador, na tabela
\texttt{ref\_faixa}. O Power BI não relaciona por chave composta, e a
saída mais óbvia seria deixar essa tabela solta e filtrar as três colunas
dentro de cada medida. O registro de decisão (ADR) do projeto escolheu
outra: \texttt{ref\_faixa} se liga a \texttt{dim\_ano} e a
\texttt{dim\_indicador} por coluna simples e vira uma segunda tabela fato,
em outro grão, que compartilha as dimensões com \texttt{fato\_indicador}.
A Faixa de porte do Município selecionado chega por \texttt{TREATAS}. Assim
o ano e o Indicador do contexto filtram sozinhos, as medidas ficam curtas e
não é preciso uma \texttt{dim\_faixa} em floco de neve, que pediria relação
bidirecional para o filtro do município subir até ela.

\begin{figure}[H]
  \centering
  \def\capturaArquivo{../powerbi/capturas/05-modelo.png}\captura
  \caption{A visão de modelo do Power BI Desktop: \NumTabelasFato{} tabelas
    fato no centro e \NumDimensoes{} dimensões conformes, em
    \NumRelacoes{} relações um-para-muitos de direção única. A referência da
    faixa compartilha o ano e o Indicador com a fato principal.}
  \label{fig:modelo}
\end{figure}

% ══════════════════════════════════════════════════════════════════════
\section{\NumMedidas{} medidas DAX, em \NumPastasMedidas{} pastas de exibição}
\label{sec:medidas}

Todo número de visual é uma medida: o modelo desencoraja a medida
implícita. As medidas ficam numa tabela própria, organizadas em pastas de
exibição e com descrição (Tabela~\ref{tab:pastas}).

\begin{table}[H]
  \centering
  \caption{As medidas por pasta de exibição, com as primeiras de cada uma.}
  \label{tab:pastas}
  \small
  \begin{tabular}{@{}lrp{9cm}@{}}
    \toprule
    Pasta & Medidas & Exemplos \\
    \midrule
    \linhasde{gerados/pastas_medidas.tex}
    \bottomrule
  \end{tabular}
\end{table}

A \textbf{Mediana da faixa} é a medida que a decisão sobre
\texttt{ref\_faixa} encurta. O ano e o Indicador já chegam pelas relações;
só a Faixa de porte, que mora em \texttt{dim\_municipio}, é levada a
\texttt{ref\_faixa} por \texttt{TREATAS}. A guarda com
\texttt{HASONEVALUE} devolve vazio, e não um \texttt{MAX} sem sentido,
quando o contexto tem mais de um ano, Indicador ou faixa.

\begin{codigo}% GERADO por scripts/numeros.py. Não edite à mão: rode o script de novo.
\begin{verbatim}
Mediana da faixa =
  IF (
    HASONEVALUE ( dim_municipio[faixa_porte] )
      && HASONEVALUE ( dim_ano[ano] )
      && HASONEVALUE ( dim_indicador[indicador_id] ),
    CALCULATE (
      MAX ( ref_faixa[mediana] ),
      TREATAS ( VALUES ( dim_municipio[faixa_porte] ), ref_faixa[faixa_porte] )
    )
  )
\end{verbatim}
\end{codigo}

A \textbf{Posição na faixa} é a fração dos outros Municípios parecidos com
valor menor que o do Município selecionado: do menor valor da faixa ao
maior. É uma posição, e não uma nota. A medida tira o filtro do município,
recoloca o da faixa e conta quantos têm valor menor, dividindo pelo tamanho
do grupo menos um, a mesma conta do \texttt{PERCENT\_RANK} do SQL. O
\texttt{COUNTROWS} de uma tabela vazia devolve \texttt{BLANK}, e não zero:
sem a soma de zero no fim, o menor valor da faixa ficava sem posição. A
conferência no motor do Power BI acusou esse defeito, que foi corrigido.

\begin{codigo}% GERADO por scripts/numeros.py. Não edite à mão: rode o script de novo.
\begin{verbatim}
Posição na faixa =
  VAR v = [Valor]
  VAR faixa = SELECTEDVALUE ( dim_municipio[faixa_porte] )
  VAR grupo =
    CALCULATETABLE (
      fato_indicador,
      REMOVEFILTERS ( dim_municipio ),
      dim_municipio[faixa_porte] = faixa
    )
  VAR n = COUNTROWS ( grupo )
  VAR menores = COUNTROWS ( FILTER ( grupo, fato_indicador[valor] < v ) ) + 0
  RETURN
    IF (
      NOT ISBLANK ( v ) && NOT ISBLANK ( faixa )
        && HASONEVALUE ( dim_ano[ano] )
        && HASONEVALUE ( dim_indicador[indicador_id] )
        && n > 1,
      DIVIDE ( menores, n - 1 )
    )
\end{verbatim}
\end{codigo}

% ══════════════════════════════════════════════════════════════════════
\section{\NumDivergenciasCompletas{} divergências em \NumCelulasCompletas{} células conferidas}
\label{sec:conferencia}

Cada número do projeto foi conferido por uma segunda implementação, ou por
um motor de verdade, comparada com a fonte, e o resultado ficou gravado no
repositório (Tabela~\ref{tab:conferencias}). As medidas DAX foram
conferidas no próprio motor do Power BI, com o Desktop aberto: primeiro
contra uma tabela de valores esperados, escrita antes de as medidas
existirem, e depois contra o SQLite, em todas as linhas da fato.

\begin{table}[H]
  \centering
  \caption{As conferências e o resultado de cada uma.}
  \label{tab:conferencias}
  \small
  \begin{tabular}{@{}p{3.7cm}p{5.4cm}rr@{}}
    \toprule
    Conferência & O que compara & Células & Resultado \\
    \midrule
    \linhasde{gerados/conferencias.tex}
    \bottomrule
  \end{tabular}
\end{table}

\begin{destaque}[As medidas devolvem os números certos]
  Na tabela de conferência, batem \NumCelulasBatem{} de
  \NumCelulasConferencia{} células, de \NumMunicipiosConferencia{}
  municípios sorteados. Contra o SQLite, a Mediana da faixa, a Posição na
  faixa e o Valor foram calculados para cada uma das \NumLinhasFato{}
  linhas da fato: \NumCelulasCompletas{} células, com
  \NumDivergenciasCompletas{} divergências.
\end{destaque}

A planilha foi conferida da mesma forma: o Excel para a web calculou as
fórmulas, e o conferidor leu os valores calculados de uma cópia por
município e os comparou com o SQLite (Seção~\ref{sec:planilha}).

% ══════════════════════════════════════════════════════════════════════
\section{O painel: \NumPaginasPainel{} páginas, com visuais definidos em código}
\label{sec:painel}

O painel abre no Power BI Desktop, a partir do projeto PBIP versionado. Os
\textbf{visuais são definidos em código (PBIR) sobre modelo em TMDL}: um
script em Python escreve cada página e cada visual, com os campos, as
medidas, o formato, o texto alternativo e a ordem de tabulação, e um
roteiro legível descreve o que cada página mostra. Um ajuste se faz no
gerador e se roda de novo; o diff mostra o que mudou em cada visual. O
Desktop serve para abrir o projeto, atualizar os dados, conferir as
medidas e exportar as capturas.

As cores seguem o kit: o Município selecionado sai no vermelho de
destaque, só em marca, e a mediana e os quartis no cinza neutro. Nenhum
visual usa cor de bom ou ruim, porque nenhum Indicador tem um lado melhor
que valha para todo município, e nenhum visual ordena municípios por nome.

\begin{figure}[H]
  \centering
  \def\capturaArquivo{../powerbi/capturas/01-municipio.png}\captura
  \caption{Página Município: \NumExemplo{} está na faixa
    \NumExemploFaixa{} habitantes, com \NumExemploNaFaixa{} municípios. Os
    cartões trazem o último valor de quatro Indicadores financeiros ao lado
    da Mediana da faixa; a linha mostra a série contra a mediana e os
    quartis.}
  \label{fig:cap-municipio}
\end{figure}

\begin{figure}[H]
  \centering
  \def\capturaArquivo{../powerbi/capturas/02-comparacao.png}\captura
  \caption{Página Comparação na faixa: \NumExemplo{} em vermelho entre os
    \NumExemploNaFaixa{} Municípios parecidos, na despesa com saúde por
    habitante. Mostra a posição no grupo, e não uma lista do pior ao
    melhor.}
  \label{fig:cap-comparacao}
\end{figure}

\begin{figure}[H]
  \centering
  \def\capturaArquivo{../powerbi/capturas/03-previsao.png}\captura
  \caption{Página Previsão da receita: R\$ \NumExemploPrevista{} milhões
    para \NumExemplo{} em \NumAnoPrevFinal, com os Intervalos de previsão de
    \NumNivelOitenta\% e de \NumNivelNoventaCinco\% e o aviso do que o
    método não conhece.}
  \label{fig:cap-previsao}
\end{figure}

\begin{figure}[H]
  \centering
  \def\capturaArquivo{../powerbi/capturas/04-como-ler.png}\captura
  \caption{Página Como ler, fontes e limites: a definição escrita de cada
    Indicador e o Teste contra o passado, em que o intervalo de
    \NumNivelOitenta\% cobriu \NumCobOitentaPior\% na origem
    \NumCobOitentaPiorOrigem.}
  \label{fig:cap-como-ler}
\end{figure}

% ══════════════════════════════════════════════════════════════════════
\section{A planilha: \NumCelulasExcelBatem{} de \NumCelulasExcel{} células iguais ao SQLite}
\label{sec:planilha}

A planilha é \textbf{gerada por script em Python e completada no Excel
para a web (tabela dinâmica e segmentação), com passo a passo para refazer
à mão}. O Excel disponível é o da web, que não grava Power Query nem roda
automação; por isso os CSV entram como tabelas nomeadas (\texttt{tFato},
\texttt{tReferencia}, \texttt{tMunicipios}, \ldots), e a fórmula lê
\texttt{tFato[valor]}, e não um intervalo fixo. O Power Query do projeto
está onde ele é natural: nas consultas do modelo do Power BI.

Na aba Painel, o leitor escolhe o município numa lista, e o PROCX traz o
resto. O valor de um município num ano e num Indicador é uma busca por
chave composta. Cada comparação entre parênteses devolve uma lista de
verdadeiros e falsos, uma por linha da tabela; o produto das três vale um
só na linha em que as três valem, e o PROCX procura esse um. É o PROCV de
chave composta sem coluna auxiliar:

\begin{codigo}% GERADO por scripts/numeros.py. Não edite à mão: rode o script de novo.
\begin{verbatim}
=PROCX(1;
       (tFato[cod_ibge]=$C$5)*(tFato[ano]=$D13)*(tFato[indicador_id]=$A13);
       tFato[valor];
       "n.d.")
\end{verbatim}
\end{codigo}

O LET dá nome a um resultado intermediário para não repetir a conta. Aqui
ele guarda o último ano com dado do município no Indicador, que as outras
colunas da linha leem:

\begin{codigo}% GERADO por scripts/numeros.py. Não edite à mão: rode o script de novo.
\begin{verbatim}
=LET(a; MÁXIMOSES(tFato[ano]; tFato[cod_ibge]; $C$5; tFato[indicador_id]; $A13);
     SE(a=0; "n.d."; a))
\end{verbatim}
\end{codigo}

A Posição na faixa da planilha faz a mesma conta da medida DAX, com LET e
SOMARPRODUTO. Na aba Explorar, uma tabela dinâmica sobre \texttt{tFato},
com segmentação por Faixa de porte, foi inserida à mão no Excel para a web.

\begin{destaque}[A planilha confere com a base]
  O Excel para a web calculou a aba Painel para \NumMunicipiosExcel{}
  municípios, e o conferidor comparou cada célula com o SQLite:
  \NumCelulasExcelBatem{} de \NumCelulasExcel{} iguais.
\end{destaque}

% ══════════════════════════════════════════════════════════════════════
\section{A tendência dos \NumJanelaEscolhido{} últimos anos erra \NumErroEscolhidoTres\% a \NumHorizontes{} anos}
\label{sec:previsao}

A Previsão da receita é a receita corrente total esperada para os
\NumHorizontes{} anos seguintes ao último observado, em R\$ milhões de
\NumAnoBaseReais, para os \NumMunicipiosComPrevisao{} municípios com
série suficiente. Foi escolhida a receita total, e não a por habitante,
porque a população muda de base no Censo 2022 e a receita total não tem
esse degrau. Concorreram \NumMetodos{} métodos simples, todos com a série
do próprio município, e cada um foi testado prevendo o passado: a previsão
parte de cada ano de \NumOrigemInicial{} a \NumOrigemFinal{} e é comparada
ao que veio depois (Tabela~\ref{tab:metodos}).

\begin{table}[H]
  \centering
  \caption{Erro médio absoluto e viés no Teste contra o passado, por
    horizonte. Repetir o último valor erra por baixo porque a receita
    cresceu no período; a tendência corrige isso.}
  \label{tab:metodos}
  \small
  \begin{tabelaplena}
  \begin{tabular}{@{}lrrrr@{}}
    \toprule
    Método & A um ano & A dois anos & A três anos & Viés a três anos \\
    \midrule
    \linhasde{gerados/teste_metodos.tex}
    \bottomrule
  \end{tabular}
  \end{tabelaplena}
\end{table}

Os Intervalos de previsão não saem de uma fórmula de erro normal: saem dos
erros que o método cometeu ao prever o passado, por horizonte e por Grupo
de porte. A receita dos municípios se move junta (inflação, ICMS,
transferências), então as \NumOrigensTeste{} origens do teste são, na
prática, \NumOrigensTeste{} choques, e não centenas de observações
independentes. Por isso o intervalo foi calibrado deixando uma origem de
fora, conferido nela e alargado até a cobertura chegar ao nível prometido
(Tabela~\ref{tab:cobertura}).

\begin{table}[H]
  \centering
  \caption{Cobertura dos Intervalos de previsão no Teste contra o passado,
    antes e depois do alargamento.}
  \label{tab:cobertura}
  \small
  \begin{tabular}{@{}lrrrl@{}}
    \toprule
    Intervalo & Sem ajuste & Fator & Com ajuste & Pior origem deixada de fora \\
    \midrule
    \linhasde{gerados/cobertura.tex}
    \bottomrule
  \end{tabular}
\end{table}

\begin{destaque}[Use o intervalo de \NumNivelNoventaCinco\% como a leitura prudente]
  Na média, o intervalo de \NumNivelOitenta\% cobre \NumCobOitentaCom\%.
  Na origem de \NumCobOitentaPiorOrigem, que vem depois de um salto de
  receita, cobriu só \NumCobOitentaPior\%. O de \NumNivelNoventaCinco\%
  cobriu \NumCobNoventaCincoPior\% na pior origem.
\end{destaque}

\begin{figure}[H]
  \centering
  \imagemplena[0.86]{03-previsao.png}
  \caption{A receita de \NumExemplo{} deve chegar a R\$
    \NumExemploPrevista{} milhões em \NumAnoPrevFinal, entre R\$
    \NumExemploInfNoventaCinco{} e \NumExemploSupNoventaCinco{} milhões no
    intervalo de \NumNivelNoventaCinco\% (entre \NumExemploInfOitenta{} e
    \NumExemploSupOitenta{} no de \NumNivelOitenta\%). Em
    \NumAnoFinal{} foram R\$ \NumExemploReceitaUltima{} milhões.}
  \label{fig:previsao}
\end{figure}

\begin{figure}[H]
  \centering
  \imagemplena[0.86]{04-teste-do-modelo.png}
  \caption{A tendência dos \NumJanelaEscolhido{} últimos anos erra
    \NumErroEscolhidoTres\% a \NumHorizontes{} anos, contra
    \NumErroIngenuoTres\% de repetir o último valor. À direita, a cobertura
    do intervalo de \NumNivelOitenta\% em cada origem deixada de fora:
    \NumCobOitentaPior\% em \NumCobOitentaPiorOrigem.}
  \label{fig:teste}
\end{figure}

% ══════════════════════════════════════════════════════════════════════
\section{O que os dados não autorizam afirmar}
\label{sec:limitacoes}

\begin{itemize}[leftmargin=*]
  \item \textbf{Municípios parecidos é só porte.} Dois municípios da mesma
    Faixa de porte podem ser muito diferentes (turismo, indústria, região
    metropolitana). A comparação diz onde o município está no grupo, e não
    por quê.
  \item \textbf{Degrau do Censo 2022.} A população anterior projeta o
    censo anterior, e a de depois vem do Censo 2022. Os Indicadores por
    habitante saltam na troca de base nos municípios que o Censo corrigiu;
    a receita total não tem o degrau.
  \item \textbf{Ideb bienal.} O Ideb sai a cada dois anos, e município sem
    rede municipal fica sem Ideb. O PIB dos municípios e o pré-natal
    terminam antes de \NumAnoFinal.
  \item \textbf{Cobertura do Intervalo de previsão.} Com \NumOrigensTeste{}
    origens, o ajuste é grosseiro: na origem de \NumCobOitentaPiorOrigem,
    o intervalo de \NumNivelOitenta\% cobriu \NumCobOitentaPior\%. O método
    não conhece mudança de lei, de alíquota, de convênio nem de gestão.
  \item \textbf{Acima da mediana não é melhor.} Gastar mais com saúde que
    a mediana não é, por si, bom nem ruim; por isso nenhum visual tem cor
    de bom ou ruim.
\end{itemize}

\end{document}
